\documentclass[a4paper,11pt]{article}
\usepackage[utf8]{inputenc}
\usepackage[french]{babel}
\usepackage[T1]{fontenc}
\usepackage{amsmath,amssymb,amsthm}
\usepackage{geometry}
\geometry{margin=2cm}
\usepackage{enumitem}
\usepackage{hyperref}
\usepackage{booktabs}
\usepackage{array}
\usepackage{xcolor}
\usepackage{listings}
\usepackage{titlesec}
\usepackage{fancyhdr}
\usepackage{lastpage}
\usepackage{graphicx}
\usepackage{etoolbox}
\AtBeginDocument{\shorthandoff{;:!?}}

\titleformat{\section}{\large\bfseries\color{blue!60!black}}{}{0pt}{}[\vspace{-0.5ex}\rule{\textwidth}{0.4pt}]
\titleformat{\subsection}{\bfseries\color{blue!40!black}}{}{0pt}{}
\titleformat{\subsubsection}{\itshape\color{blue!30!black}}{}{0pt}{}

\lstdefinestyle{monstyle}{
  frame=single,
  numbers=left,
  numbersep=5pt,
  breaklines=true,
  basicstyle=\small\ttfamily,
  backgroundcolor=\color{gray!5},
  rulecolor=\color{gray!40},
  columns=fullflexible,
  keepspaces=true,
  showstringspaces=false
}
\lstset{style=monstyle}

\pagestyle{fancy}
\fancyhf{}
\fancyhead[L]{\small STA211 -- Recherche de r\`egles d'association}
\fancyhead[R]{\small Fiche r\'ecapitulative}
\fancyfoot[C]{\thepage/\pageref{LastPage}}
\renewcommand{\headrulewidth}{0.4pt}

\theoremstyle{definition}
\newtheorem*{definition}{D\'efinition}
\theoremstyle{remark}
\newtheorem*{remark}{Remarque}
\newtheorem*{important}{\`A retenir}

\title{\textbf{Fiche r\'ecapitulative -- STA211}\\[4pt]
\large Recherche de r\`egles d'association}
\author{Vincent Audigier, Nd\`eye Niang}
\date{12 juin 2026}

\begin{document}
\maketitle
\thispagestyle{empty}

\begin{abstract}
Cette fiche pr\'esente les m\'ethodes de recherche de r\`egles
d'association, initialement d\'evelopp\'ees pour l'analyse des
paniers d'achat (Agrawal et al., 1993). Elle couvre les concepts
fondamentaux (items, transactions, itemsets, support, confiance),
l'algorithme Apriori et sa propri\'et\'e d'anti-monotonie, les
algorithmes alternatifs (représentation horizontale et verticale),
les indices de pertinence (lift, conviction, leverage) ainsi que des
exemples d'application sur donn\'ees supermarch\'e et German Credit.
\end{abstract}

\vfill
\tableofcontents
\newpage

% ===================================================================
\section{Introduction}
% ===================================================================

Les m\'ethodes de recherche de r\`egles d'association ont \'et\'e
propos\'ees par \cite{Agrawal1993} pour d\'ecouvrir quels produits
\'etaient achet\'es conjointement dans les bases de donn\'ees des
ventes de supermarch\'e. Elles permettent d'extraire des r\`egles
du type~:

\begin{quote}
  \textit{``Lorsqu'un client ach\`ete du pain et du beurre, alors
    neuf fois sur dix il ach\`ete en m\^eme temps du lait.''}
\end{quote}

Bien que d\'evelopp\'ees initialement pour le marketing, ces
m\'ethodes s'appliquent \'egalement au webmining (analyse de
consultation de pages web), au text mining (cooccurrences de
termes), etc.

% ===================================================================
\section{D\'efinitions et concepts de base}
% ===================================================================

\subsection{Items, transactions et r\`egles}

On d\'efinit $I = \{it_1, it_2, \dots, it_p\}$ l'ensemble des
\textbf{items} (produits d'un supermarch\'e par exemple). Une
\textbf{transaction} $tr_j$ est un sous-ensemble non vide de $I$~;
on note $T = \{tr_1, tr_2, \dots, tr_n\}$ l'ensemble des
$n$ transactions (contenus de caddies).

Une \textbf{r\`egle d'association} est une implication
$A \rightarrow B$ o\`u $A \subset I$, $B \subset I$ et
$A \cap B = \emptyset$. $A$ est la pr\'emisse (ant\'ec\'edent),
$B$ la conclusion (cons\'equent). Les ensembles $A$ et $B$ sont
appel\'es \textbf{itemsets}.

\subsection{Support et confiance}

\begin{definition}[Support]
  Le support d'un itemset $X$ est la proportion de transactions
  contenant $X$~:
  \begin{equation}
    \text{supp}(X) = \frac{|\{t \in T : X \subseteq t\}|}{|T|}
  \end{equation}
\end{definition}

\begin{definition}[Confiance]
  La confiance d'une r\`egle $A \rightarrow B$ est la proportion
  de transactions contenant $B$ parmi celles contenant $A$~:
  \begin{equation}
    \text{conf}(A \rightarrow B) = \frac{\text{supp}(A \cup B)}{\text{supp}(A)}
  \end{equation}
\end{definition}

\begin{important}
  Le support mesure la \textbf{g\'en\'eralit\'e} d'une r\`egle~;
  la confiance mesure sa \textbf{fiabilit\'e}. On impose
  g\'en\'eralement des seuils minimaux $\text{supp}_{\min}$ et
  $\text{conf}_{\min}$.
\end{important}

% ===================================================================
\section{L'algorithme Apriori}
% ===================================================================

L'algorithme Apriori (\cite{Agrawal1994}) proc\`ede en deux
\'etapes~: g\'en\'eration des itemsets fr\'equents puis extraction
des r\`egles.

\subsection{Recherche des itemsets fr\'equents}

On explore le treillis des itemsets par parcours en largeur (de
taille 1 jusqu'\`a $p$). \`A chaque niveau $k$, on g\'en\`ere les
candidats de taille $k$ \`a partir des itemsets fr\'equents de
taille $k-1$, puis on \'elague ceux dont le support est inf\'erieur
\`a $\text{supp}_{\min}$.

La propri\'et\'e d'\textbf{anti-monotonie} garantit l'\'elagage~:
\begin{important}
  Tout sur-ensemble d'un itemset non fr\'equent est n\'ecessairement
  non fr\'equent.
\end{important}

Si un itemset de taille $k-1$ a un support $< \text{supp}_{\min}$,
tous ses sur-ensembles de taille $k$ peuvent \^etre \'elimin\'es
sans calcul suppl\'ementaire.

\subsection{Extraction des r\`egles}

\`A partir des itemsets fr\'equents, on g\'en\`ere toutes les
r\`egles $A \rightarrow B$ telles que
$\text{conf}(A \rightarrow B) \geq \text{conf}_{\min}$.

% ===================================================================
\section{Algorithmes alternatifs}
% ===================================================================

\subsection{Repr\'esentation horizontale}

Les donn\'ees sont stock\'ees sous forme d'une liste de transactions
(une ligne par transaction, chaque ligne contient la liste des items
achet\'es). Apriori utilise cette repr\'esentation.

\subsection{Repr\'esentation verticale}

Les donn\'ees sont stock\'ees sous forme d'une liste d'items~: pour
chaque item, on stocke la liste des identifiants de transactions
qui le contiennent (format \emph{tid-list}). Avantage~: le support
d'un itemset s'obtient par intersection des tid-lists.

% ===================================================================
\section{Indices de pertinence des r\`egles}
% ===================================================================

\subsection{Faiblesse de l'approche support-confiance}

Une r\`egle peut avoir une confiance \'elev\'ee mais \^etre
trompeuse si le cons\'equent est fr\'equent dans la population.
Par exemple, si 90\,\% des clients ach\`etent du lait, une r\`egle
$X \rightarrow \text{lait}$ aura une confiance d'au moins 90\,\%
quel que soit $X$.

\subsection{Indices compl\'ementaires}

\begin{definition}[Lift]
  Le lift mesure l'\'ecart \`a l'ind\'ependance~:
  \begin{equation}
    \text{lift}(A \rightarrow B) = \frac{\text{supp}(A \cup B)}{%
      \text{supp}(A)\,\text{supp}(B)}
  \end{equation}
  Interpr\'etation~:
  \begin{itemize}
    \item $\text{lift} > 1$~: association positive (items achet\'es
      plus souvent ensemble que sous l'ind\'ependance).
    \item $\text{lift} = 1$~: ind\'ependance.
    \item $\text{lift} < 1$~: association n\'egative.
  \end{itemize}
\end{definition}

\begin{definition}[Conviction]
  \begin{equation}
    \text{conviction}(A \rightarrow B) = \frac{1 - \text{supp}(B)}{%
      1 - \text{conf}(A \rightarrow B)}
  \end{equation}
  Mesure la d\'ependance directionnelle~: plus la conviction est
  grande, plus la r\`egle est robuste.
\end{definition}

\begin{definition}[Leverage]
  \begin{equation}
    \text{leverage}(A \rightarrow B) = \text{supp}(A \cup B) -
    \text{supp}(A)\,\text{supp}(B)
  \end{equation}
\end{definition}

% ===================================================================
\section{Exemples}
% ===================================================================

\subsection{Supermarch\'e}

Soit 5 transactions~:

\begin{center}
\begin{tabular}{ll}
\toprule
Transaction & Items \\
\midrule
$tr_1$ & $\{ \text{pain}, \text{lait}, \text{beurre}, \text{bi\`ere} \}$ \\
$tr_2$ & $\{ \text{pain}, \text{beurre} \}$ \\
$tr_3$ & $\{ \text{bi\`ere}, \text{cacahu\`etes} \}$ \\
$tr_4$ & $\{ \text{pain}, \text{lait}, \text{beurre} \}$ \\
$tr_5$ & $\{ \text{pain}, \text{lait} \}$ \\
\bottomrule
\end{tabular}
\end{center}

\noindent Avec $\text{supp}_{\min} = 0{,}4$~:
\begin{itemize}
  \item $\text{supp}(\{\text{pain}\}) = 4/5 = 0{,}8$~: fr\'equent.
  \item $\text{supp}(\{\text{pain}, \text{beurre}\}) = 3/5 = 0{,}6$~: fr\'equent.
  \item $\text{supp}(\{\text{cacahu\`etes}\}) = 1/5 = 0{,}2$~: non fr\'equent.
\end{itemize}

R\`egle $\{\text{pain}\} \rightarrow \{\text{beurre}\}$~:
\begin{itemize}
  \item $\text{supp} = 3/5 = 0{,}6$
  \item $\text{conf} = 3/4 = 0{,}75$
  \item $\text{lift} = 0{,}6 / (0{,}8 \times 0{,}6) = 1{,}25 > 1$
\end{itemize}

\subsection{German Credit}

Les r\`egles d'association peuvent \^etre extraites apr\`es
discr\'etisation des variables quantitatives avec l'algorithme
Apriori via le package R \texttt{arules}.

\begin{lstlisting}[language=R]
library(arules)

# Discretisation des variables quantitatives
don_disc <- discretizeDF(don, default = list(method = "interval",
                         breaks = 3))

# Extraction des regles
regles <- apriori(don_disc,
                  parameter = list(supp = 0.05, conf = 0.3,
                                   target = "rules"))
inspect(head(regles, 5))
\end{lstlisting}

% ===================================================================
\section{Conclusion}
% ===================================================================

Les m\'ethodes de recherche de r\`egles d'association permettent
d'extraire des connaissances sous forme de r\`egles
interpr\'etables \`a partir de grandes bases de donn\'ees.
L'algorithme Apriori reste la m\'ethode de r\'ef\'erence gr\^ace
\`a son \'elagage par anti-monotonie. Les indices de pertinence
(lift, conviction, leverage) compl\`etent l'approche
support-confiance pour \'eviter les r\`egles trompeuses.

Applications principales~:
\begin{itemize}
  \item Analyse de paniers d'achat (marketing).
  \item Webmining (analyse de navigation).
  \item Text mining (cooccurrences de termes).
  \item Analyse de donn\'ees m\'edicales (associations de
    sympt\^omes).
\end{itemize}

% ===================================================================
\section*{Questions cl\'es (QCM)}
% ===================================================================

\addcontentsline{toc}{section}{Questions cl\'es (QCM)}

\begin{enumerate}

  \item Que repr\'esente le support d'un itemset $X$~?
    \begin{enumerate}
      \item La probabilit\'e de $X$ sachant $Y$
      \item La proportion de transactions contenant $X$
      \item La force de l'association $X \Rightarrow Y$
      \item Le nombre d'items dans $X$
    \end{enumerate}
    \textbf{R\'eponse~: b (proportion de transactions contenant $X$).}

  \item Que mesure la confiance d'une r\`egle $A \rightarrow B$~?
    \begin{enumerate}
      \item $\displaystyle \frac{\text{supp}(A \cup B)}{\text{supp}(A)}$
      \item $\displaystyle \frac{\text{supp}(A)}{\text{supp}(B)}$
      \item $\displaystyle \frac{\text{supp}(A \cup B)}{\text{supp}(B)}$
      \item $\displaystyle \frac{\text{supp}(A)}{\text{supp}(A \cup B)}$
    \end{enumerate}
    \textbf{R\'eponse~: a ($\text{supp}(A \cup B) / \text{supp}(A)$).}

  \item Quelle propri\'et\'e permet l'\'elagage dans Apriori~?
    \begin{enumerate}
      \item La monotonie du support
      \item L'anti-monotonie du support
      \item La sym\'etrie de la confiance
      \item La transitivit\'e du lift
    \end{enumerate}
    \textbf{R\'eponse~: b (anti-monotonie~: tout sur-ensemble d'un
    itemset non fr\'equent est non fr\'equent).}

  \item Un lift $> 1$ indique~:
    \begin{enumerate}
      \item Une association n\'egative
      \item Une association positive
      \item Aucune association
      \item Une erreur de calcul
    \end{enumerate}
    \textbf{R\'eponse~: b (association positive).}

  \item Quelle est la particularit\'e de la repr\'esentation
    verticale des donn\'ees~?
    \begin{enumerate}
      \item Chaque ligne est une transaction
      \item Pour chaque item, on stocke la liste des transactions
        (tid-list)
      \item Les donn\'ees sont stock\'ees sous forme matricielle
      \item Elle n\'ecessite moins de m\'emoire qu'Apriori
    \end{enumerate}
    \textbf{R\'eponse~: b (format tid-list).}

  \item Quel package R impl\'emente l'algorithme Apriori~?
    \begin{enumerate}
      \item \texttt{arules}
      \item \texttt{cluster}
      \item \texttt{FactoMineR}
      \item \texttt{MASS}
    \end{enumerate}
    \textbf{R\'eponse~: a (\texttt{arules}).}

  \item Que mesure la conviction d'une r\`egle~?
    \begin{enumerate}
      \item Le support conditionnel
      \item La d\'ependance directionnelle entre pr\'emisse et
        conclusion
      \item La fr\'equence de l'itemset
      \item La dissimilarit\'e entre items
    \end{enumerate}
    \textbf{R\'eponse~: b (d\'ependance directionnelle).}

  \item Dans l'exemple du supermarch\'e avec 5 transactions,
    quel est le support de $\{\text{pain}, \text{beurre}\}$~?
    \begin{enumerate}
      \item $3/5 = 0{,}6$
      \item $4/5 = 0{,}8$
      \item $2/5 = 0{,}4$
      \item $1/5 = 0{,}2$
    \end{enumerate}
    \textbf{R\'eponse~: a ($3/5$).}

  \item Quel est l'inconv\'enient de l'approche support-confiance~?
    \begin{enumerate}
      \item Elle est trop complexe
      \item Une r\`egle peut avoir une confiance \'elev\'ee mais
        \^etre trompeuse si le cons\'equent est fr\'equent
      \item Elle ne fonctionne que pour les donn\'ees binaires
      \item Elle n\'ecessite trop de m\'emoire
    \end{enumerate}
    \textbf{R\'eponse~: b (confiance \'elev\'ee mais r\`egle
    trompeuse).}

  \item Qui a propos\'e l'algorithme Apriori~?
    \begin{enumerate}
      \item Vigneau et Qannari
      \item Agrawal, Imieli\'nski et Swami
      \item Chavent et al.
      \item Breiman
    \end{enumerate}
    \textbf{R\'eponse~: b (Agrawal, Imieli\'nski et Swami, 1993).}

\end{enumerate}

% ===================================================================
\begin{thebibliography}{9}
% ===================================================================

\bibitem{Agrawal1993}
  R.~Agrawal, T.~Imieli\'nski, A.~Swami.
  \newblock \emph{Mining association rules between sets of items in
    large databases}.
  \newblock Proc. ACM SIGMOD, 22(2):207--216, 1993.

\bibitem{Agrawal1994}
  R.~Agrawal, R.~Srikant.
  \newblock \emph{Fast algorithms for mining association rules in
    large databases}.
  \newblock Proc. VLDB, 487--499, 1994.

\bibitem{Plasse2006}
  M.~Plasse.
  \newblock \emph{Regroupement d'items dans les probl\`emes de
    recherche de r\`egles d'association}.
  \newblock Th\`ese de doctorat, CNAM, 2006.

\end{thebibliography}

\end{document}
